\section{Results and Discussion}\label{sec:Results}
We describe our setup (\cref{sec:Setup}), compare our predictions with measurements (\cref{sec:Measurements}), and discuss the costs of subdivision (\cref{sec:Discussion}).

\subsection{Setup}\label{sec:Setup}
Our rasterizer follows \cref{lst:Interpolation} and computes in double precision.
It renders $960\times400$ pixels and compares affine with perspective-correct texture coordinates at the pixel centers.
The floor of \cref{fig:Teaser} is 2~units wide and spans the depths from 1 to 4, i.e., $r = 4$.
Its texture has $256\times256$ texels, so its diagonal has the length $\ell = 362$~texels.
The script \texttt{scripts/sample\_paper.py} renders all images and writes all data of this paper.

\subsection{Predictions and Measurements}\label{sec:Measurements}
\Cref{fig:Plots}~(a) compares \cref{eq:MaxError} with the measured maximum for floors with depth ratios from 1.25 to 16.
The measurements stay below the predictions by at most \qty{0.05}{\percent}, since the pixel centers miss the exact maximum.
The bound $(r - 1)/4$ of \cref{sec:Subdivision} is tight only for small depth ratios.

\begin{figure*}[t]
    \centering
    % Plots of Sec. 4. pgfplots reads the data that scripts/sample_paper.py writes to figs/data/,
% so the plots follow the data after every run of the script.
\begin{tikzpicture}
\begin{groupplot}[
    group style={group size=2 by 1, horizontal sep=20mm},
    scale only axis, width=0.36\textwidth, height=32mm,
    xmode=log, log basis x=2, log ticks with fixed point,
    grid=major, grid style={PaletteGray!25},
    tick align=outside, tick pos=left,
    every axis plot/.append style={line width=0.8pt},
    legend cell align=left,
    table/col sep=comma,
    % Numbers as siunitx prints them: no separator below 10000.
    /pgf/number format/1000 sep={\,}, /pgf/number format/min exponent for 1000 sep=4,
]
\nextgroupplot[
    xmin=1, xmax=16, xtick={1, 2, 4, 8, 16},
    ymin=0, ymax=0.7, ytick={0, 0.2, 0.4, 0.6},
    xlabel={Depth ratio $r$},
    ylabel={Max.\ error $\delta_\mathrm{max}$},
]
\addplot[PaletteBlue, domain=1:16, samples=200] {(sqrt(x) - 1) / (sqrt(x) + 1)};
\addplot[PaletteDarkBlue, dashed, domain=1:16, samples=200] {(x - 1) / 4};
\addplot[PaletteOrange, only marks, mark=o, mark size=2pt] table[x=ratio, y=measured] {figs/data/edge_error.csv};
\nextgroupplot[
    xmin=1, xmax=32, xtick={1, 2, 4, 8, 16, 32},
    ymode=log, ymin=0.1, ymax=1000, ytick={0.1, 1, 10, 100, 1000},
    xlabel={Pieces per edge $n$},
    ylabel={Max.\ error $E_n$ [texels]},
    legend pos=north east, legend columns=2,
]
\addplot[PaletteBlue] table[x=n, y=predicted] {figs/data/subdivision.csv};
\addplot[PaletteDarkBlue, dashed] table[x=n, y=bound] {figs/data/subdivision.csv};
\addplot[PaletteOrange, only marks, mark=o, mark size=2pt] table[x=n, y=measured] {figs/data/subdivision.csv};
\addplot[PaletteGray, dotted, domain=1:32, samples=2] {0.5};
\legend{Predicted, Bound, Measured, Half a texel}
\end{groupplot}
\node[anchor=north] at ([yshift=-1mm]group c1r1.outer south) {(a) One edge};
\node[anchor=north] at ([yshift=-1mm]group c2r1.outer south) {(b) Subdivided floor};
\end{tikzpicture}

    \caption{\captiontitle{Predicted and Measured Errors}
        (a)~The largest error along an edge as a fraction of the edge, over the depth ratio $r$ of its endpoints.
        (b)~The largest error of the floor of \cref{fig:Teaser} in texels after splitting it into $n\times n$ cells.
        Solid lines show \cref{eq:MaxError}, dashed lines the bounds of \cref{sec:Subdivision}, and circles our measurements.
    }
    \label{fig:Plots}
\end{figure*}

\Cref{tab:Subdivision} and \cref{fig:Plots}~(b) show the largest error after subdividing the floor into $n\times n$ cells.
Prediction and measurement differ by at most \qty{0.1}{\percent}.
The error falls quadratically, as \cref{eq:SubdivisionBound} predicts.
The bound overestimates the error by a factor of 2.25 for $n = 1$ and by \qty{5}{\percent} for $n = 32$.
As \cref{eq:Pieces} predicts, $24\times24$ cells keep the error below half a texel, while $16\times16$ cells do not.

\begin{table}[t]
    \centering
    % Generated by scripts/sample_paper.py; do not edit by hand.
% Modern table: \toprule, \midrule, \bottomrule, and \cmidrule to group the error columns.
\begin{tabular}{@{}rrrrr@{}}
\toprule
 & & \multicolumn{3}{c}{Max.\ error [texels]} \\
\cmidrule(l){3-5}
$n$ & Triangles & Predicted & Bound & Measured $\downarrow$ \\
\midrule
1 & \num{2} & \num{120.680} & \num{271.529} & \num{120.678} \\
2 & \num{8} & \num{40.756} & \num{67.882} & \num{40.756} \\
4 & \num{32} & \num{12.581} & \num{16.971} & \num{12.581} \\
8 & \num{128} & \num{3.595} & \num{4.243} & \num{3.595} \\
16 & \num{512} & \num{0.972} & \num{1.061} & \num{0.971} \\
24 & \num{1152} & \num{0.444} & \num{0.471} & \textbf{\num{0.444}} \\
32 & \num{2048} & \num{0.253} & \num{0.265} & \textbf{\num{0.253}} \\
\bottomrule
\end{tabular}

    \caption{\captiontitle{Error After Subdivision}
        We split the floor of \cref{fig:Teaser} into $n\times n$ cells of two triangles each.
        Predicted is \cref{eq:MaxError} for the diagonal of a cell next to the camera, bound is \cref{eq:SubdivisionBound}.
        Bold marks errors below half a texel.
    }
    \label{tab:Subdivision}
\end{table}

\subsection{Discussion}\label{sec:Discussion}
Subdivision trades triangles for accuracy.
Half a texel costs \num{1152} instead of two triangles for a single floor.
By \cref{eq:Pieces}, the number of triangles $2n^2$ grows linearly with the length $\ell$ of an edge in the texture, and thus with the texture resolution.
With $4096\times4096$ texels, \cref{eq:Pieces} asks for $94\times94$ cells, i.e., \num{17672} triangles.
Perspective-correct interpolation needs no subdivision; it costs a division per pixel, independent of the texture resolution.

The error concentrates near the camera, where the pieces have the largest depth ratios.
With $24\times24$ cells, the farthest cells of the floor err by at most 0.12~texels, about a quarter of the nearest ones.
Pieces of equal length thus waste triangles far from the camera.
In a mesh, a renderer can choose $n$ for each edge, since \cref{eq:Pieces} needs only the texture coordinates and the depths of its endpoints.
Then, distant edges stay coarse.
